\documentclass[10pt,a4paper]{amsart}
\usepackage[margin=19mm]{geometry}
\usepackage{amsmath,amssymb,mathtools}
\usepackage[scr=rsfs,cal=euler]{mathalfa}
\usepackage{xfrac}
\usepackage[unicode,hidelinks,bookmarks=false]{hyperref}
\newcommand{\ie}{i.e.\ }
\newcommand{\cf}{cf.\ }
\newcommand{\resp}{respectively\ }
\newcommand{\Subsection}{\subsection}
\providecommand{\nobreakdash}{\nobreak\hbox{-}}
\setlength{\emergencystretch}{2em}
\multlinegap0pt
\usepackage{amssymb}
\usepackage{mathtools}
\usepackage[shortlabels]{enumitem}
\newtheorem{corollary}{Corollary}
\newtheorem{lemma}{Lemma}
\newcommand{\dcup}{\mathbin{\dot\cup}}
\newcommand{\eps}{\varepsilon}
\title{A linear bound for nested cycles without geometric crossings}
\author{Jiangdong Ai, Gregory Gutin, Yiming Hao}
\date{Source: arXiv:2609.02234v2 (CC BY 4.0)}
\begin{document}
\maketitle

\noindent\emph{Source and licence.} A linear bound for nested cycles without geometric crossings. Version arXiv:2609.02234v2 (CC BY 4.0), distributed by arXiv under the Creative Commons Attribution 4.0 International licence, http://creativecommons.org/licenses/by/4.0/, as recorded in the arXiv licence metadata for this exact version and shown on the abstract page https://arxiv.org/abs/2609.02234v2. Full licence notice: \texttt{licenses/arxiv-2609.02234-CC-BY-4.0.txt}.\par\smallskip

\noindent\textit{Editorial scope.} Counted unit: the article's lemma lem:fullfan (source0.tex lines 565--571). The article's complete original proof of it is delivered with it (source0.tex lines 572--614, one \texttt{proof} environment of 43 lines). No mathematics is written, repaired or re-derived: the statement and the proof are the article's own lines, in the article's own notation, and no retained line is substituted or re-flowed.  Retained uncounted as the source-local context the proof needs: lines 370--372; lines 404--406; lines 457--461; lines 484--486; lines 489--496; lines 491--495; lines 519--524; lines 547--550; lines 554--558.  Nothing was omitted from any retained span, so the package carries no omission record.  No retained line contains a citation, so the wrapper carries no bibliography and no bibliography entry.  No retained line contains a figure environment, an image-inclusion command or drawing code, so no image is delivered and the package declares no asset dependency.  Editorial formatting only, in addition: an AMS \texttt{amsart} preamble that restates the article's own declarations and macros used by the retained lines, namely the environments \texttt{corollary}, \texttt{lemma} on separate counters, together with the standard packages those lines need.  The retained environments print in this document as Corollary 1, Lemma 1 in order of appearance, whereas the article prints the same items with the numbering of its own full document.  The complete native source of the article as distributed in the arXiv e-print for this exact version is bundled unchanged as \texttt{sources/209184/source0.tex}.
\begin{equation}\label{eq:LS}
e(F[W])\le\lambda|W|\qquad(0<|W|\le\tau_j(N)),
\end{equation}

\begin{corollary}\label{cor:reservedlocal}
Fix $j\ge1$ and $r\ge32$, and let $\lambda=8r$ and $K=128M_j\lambda$. Suppose that $F\subseteq H$, $|F|\le N$, $N\ge N_j$, condition~\eqref{eq:LS} holds, and $e(F)>K|F|$. Then $F$ contains a $j$-layer family with outermost cycle of length at most $P_j(N)$, and there is a set $B\subseteq V(F)$ disjoint from its outer vertex set $X$ such that $d_H(x,B)\ge2r$ for every $x\in X$ and $\delta(H[B])\ge r$.
\end{corollary}

\begin{equation}\label{eq:indhyp}
e(F)>\kappa_j|F|
\quad\Longrightarrow\quad
F\text{ contains }j\text{ nested cycles without geometric crossings}.
\end{equation}

\begin{equation}\label{eq:MIN}
e(H[S])\le\lambda|S|\qquad(0<|S|<q).
\end{equation}

\begin{lemma}\label{lem:peelX}
Starting from $H-X$, repeatedly delete a vertex of current degree less than $r$. Let $E\subseteq V(H)\setminus X$ be the set of additionally deleted vertices and let $R=V(H)\setminus(X\cup E)$ be the final remaining set. Thus $V(H)=X\dcup E\dcup R$, and
\begin{equation}\label{eq:CORE}
|E|<q,\qquad
d_H(x,R)\ge2r\quad(x\in X),\qquad
\delta(H[R])\ge r.
\end{equation}
\end{lemma}

\begin{equation}
|E|<q,\qquad
d_H(x,R)\ge2r\quad(x\in X),\qquad
\delta(H[R])\ge r.
\end{equation}

\begin{lemma}\label{lem:smallsets}
Every nonempty $S\subseteq V(H)$ with $2|S|<q$ satisfies
\begin{equation}\label{eq:smallMIN}
e(H[S])\le256\kappa_j|S|.
\end{equation}
\end{lemma}

\begin{equation}\label{eq:INV}
|U|\le|H|/2,\quad |N_H(U)|\le s
\quad\Longrightarrow\quad |U|<F_\eps(s).
\end{equation}

\begin{equation}\label{eq:q0}
\mathcal N(q)\ge N_j,\qquad
\tau_j(\mathcal N(q))<q,\qquad
P_j(\mathcal N(q))<q.
\end{equation}

\begin{lemma}\label{lem:fullfan}
Assume~\eqref{eq:indhyp} and the choices above. Let $H$ be a weak expander on $n$ vertices with $\delta(H)\ge D$, and let $X,E,R$ come from a minimum eligible family and Lemma~\ref{lem:peelX}; write $q=|X|$. Every set $T\subseteq R$ with
\begin{equation}\label{eq:target}
|T|>F_\eps(4q)+2q
\end{equation}
admits two $X$--$T$ paths from each vertex of $X$, such that all $2q$ paths lie in $H[X\cup R]$, are disjoint outside $X$, have distinct endpoints in $T$, and have no other vertex of $X$ in their interiors.
\end{lemma}

\begin{proof}
Construct a directed network with source $\sigma$ and sink $\omega$; an arc $uv$ is directed from $u$ to $v$. Each arc $\sigma x$, for $x\in X$, has capacity two. Replace each $v\in R$ by two distinct network vertices, its in-copy $v^-$ and out-copy $v^+$, with an arc $v^-v^+$ of capacity one. Include arcs $xv^-$ for each edge $xv$ from $X$ to $R$, both arcs $u^+v^-,v^+u^-$ for each edge $uv$ of $H[R]$, and arcs $t^+\omega$ for $t\in T$. Give all these remaining arcs capacity $2q+1$. An integral flow assigns integer arc values between zero and capacity, with conservation except at $\sigma,\omega$; its value is the net outflow from $\sigma$. A cut is a partition separating $\sigma$ from $\omega$; its capacity sums the capacities of arcs directed from the source side to the sink side.

An integral flow of value $2q$ gives the required paths after discarding flow cycles and, if necessary, stopping each path at its first target vertex. Suppose that no such flow exists, and take a minimum cut of capacity less than $2q$. No arc of capacity $2q+1$ crosses the cut. We may also assume that whenever $v^+$ is on the source side, $v^-$ is on that side: moving $v^-$ there cannot increase the cut capacity, since its only outgoing arc is the split arc to $v^+$.

Let $U\subseteq X$ be the roots on the source side, $B\subseteq R$ the vertices with both copies there, and $Z\subseteq R$ the vertices with only their in-copy there. The cut capacity is $2(q-|U|)+|Z|<2q$. Hence
\begin{equation}\label{eq:cut}
\begin{gathered}
U\ne\varnothing,\qquad |Z|<2|U|,\\
\Gamma_H(U)\cap R\subseteq B\cup Z,\qquad
N_{H[R]}(B)\subseteq Z,\qquad B\cap T=\varnothing.
\end{gathered}
\end{equation}
The neighbourhood assertions follow because no large-capacity arc crosses the cut.

In $H$, the external neighbourhood of $B$ is contained in $X\cup E\cup Z$, which has fewer than $4q$ vertices. If $|B|>n/2$, then $B'=R\setminus(B\cup Z)$ has fewer than $n/2$ vertices. Since $N_{H[R]}(B)\subseteq Z$, no vertex of $B'$ is adjacent to $B$, so $N_H(B')\subseteq X\cup E\cup Z$ as well. Moreover, $B'$ contains $T\setminus Z$, whose order exceeds $F_\eps(4q)$ by~\eqref{eq:target}. This contradicts~\eqref{eq:INV}. Therefore $|B|\le n/2$, and~\eqref{eq:INV} gives $|B|<F_\eps(4q)$. Let $t=|U|+|B|$.

\medskip
\noindent\emph{Case 1: $t<q/8$.}
Let $S=U\cup B\cup Z$. Then $0<|S|<3t<3q/8$, so $2|S|<q$. By~\eqref{eq:CORE} and~\eqref{eq:cut}, at least $2r|U|$ edges join $U$ to $B\cup Z$. Also, the edges inside $H[B\cup Z]$ incident with $B$ number at least $r|B|/2$. These edge collections are disjoint. Since $|Z|<2|U|$,
\[
e(H[S])\ge2r|U|+\frac r2|B|
>\frac r2|S|\ge512\kappa_j|S|,
\]
contrary to Lemma~\ref{lem:smallsets}.

\medskip
\noindent\emph{Case 2: $t\ge q/8$.}
Let $W=X\cup E\cup B\cup Z$. All neighbours in $H$ of every vertex of $U\cup B$ lie in $W$. The original minimum degree therefore gives
\[
2e(H[W])\ge Dt,\qquad
|W|<2q+2t,\qquad |W|<\mathcal N(q).
\]
The last bound uses $|B|<F_\eps(4q)$, $|E|<q$ and $|Z|<2q$.

If $q\le q_0$, a vertex of the nonempty set $U$ would have all its neighbours in a set of fewer than $\mathcal N(q_0)<D$ vertices, a contradiction. Thus $q>q_0$. Moreover, the preceding bounds give
\[
\frac{e(H[W])}{|W|}>\frac{Dt}{4(q+t)}\ge\frac D{36}>K.
\]
At ambient scale $N=\mathcal N(q)$, conditions~\eqref{eq:q0} and~\eqref{eq:MIN} give every local edge bound required by Corollary~\ref{cor:reservedlocal}. That corollary produces an eligible $j$-layer family in $H[W]$ whose outermost cycle has length at most $P_j(\mathcal N(q))<q$, again contradicting minimality.

Both cases are impossible. The maximum-flow min-cut theorem yields a flow of value $2q$. The resulting paths lie in $H[X\cup R]$, since the network contains no vertex of $E$.
\end{proof}

\end{document}
